\documentclass{resume}

\usepackage{hyperref}
\usepackage{xcolor}
\urlstyle{same}
\hypersetup{
  colorlinks=true,
  breaklinks=true,
  linkcolor={orange!60!black},
  citecolor={orange!60!black},
  urlcolor={orange!60!black}
}


\definecolor{grey}{rgb}{0.5, 0.5, 0.5}

\usepackage{tgadventor}
\renewcommand{\familydefault}{\sfdefault}

\author{ Paul McCarthy }

\address{


\textcolor{grey}{10 Newton Street}\\

\textcolor{grey}{Llanberis, LL55 4HN}\\

\textcolor{grey}{Wales, United Kingdom}\\

\textcolor{grey}{+44 7474795710}\\

}
{

\textcolor{grey}{DOB 26th July 1982}\\

\textcolor{grey}{Australian citizen}\\

\textcolor{grey}{UK permanent resident}\\


\mbox{\small\url{https://github.com/pauldmccarthy/}}\\

\mbox{\small\url{https://git.fmrib.ox.ac.uk/paulmc/}}\\

\href{mailto:pauldmccarthy@gmail.com}{\small\nolinkurl{pauldmccarthy@gmail.com}}\\

\href{mailto:paul.mccarthy@ndcn.ox.ac.uk}{\small\nolinkurl{paul.mccarthy@ndcn.ox.ac.uk}}\\

}

\begin{document}

\maketitle

\begin{category}{Summary}
\citemnobullet I am a research software engineer with a wide range of experience in
different domains and programming languages. I am presently employed
at the \href{https://www.oxcin.ox.ac.uk}{Oxford Centre for Integrative Neuroimaging
(OxCIN/FMRIB)}, where I work on a range of projects related to
the analysis and visualisation of functional and structural MRI
data. I have previously held roles in modelling and simulation,
microcontroller programming, web development, and data analysis.
\end{category}


\begin{category}{Timeline}
\citemnobullet

%% skip last section ("outside of work")


\item[] \textbf{2014-present - Research software engineer}\newline
\textcolor{grey}{\textit{Oxford Centre for Integrative Neuroimaging (OxCIN/FMRIB), UK}}



\item[] \textbf{2010-2013 - PhD}\newline
\textcolor{grey}{\textit{Department of Computer Science, University of Otago, Dunedin, New Zealand}}



\item[] \textbf{2007-2010 - Software engineer}\newline
\textcolor{grey}{\textit{CSIRO, Hobart, Australia}}



\item[] \textbf{2006-2007 - Software engineer}\newline
\textcolor{grey}{\textit{Australian Defence Force Academy, Canberra, Australia}}



\item[] \textbf{2002-2006 - Master of Science}\newline
\textcolor{grey}{\textit{Department of Computer Science, University of South Australia, Adelaide, Australia}}




\end{category}



\begin{category}{2014-present - Research software engineer}
\citemnobullet
\begin{changemargin}{-0.5in}{0.5in}
\textcolor{grey}{\textit{Oxford Centre for Integrative Neuroimaging (OxCIN/FMRIB), UK}}
\end{changemargin}

\citemnobullet My current position involves software development, management, and
maintenance on a range of projects related to \href{https://fsl.fmrib.ox.ac.uk/fsl/docs/}{FSL},
the FMRIB Software Library.\newline

Some of my software contributions to FMRIB and FSL have been:\newline

\begin{itemize}
\item \href{https://open.oxcin.ox.ac.uk/pages/fsl/fsleyes/fsleyes/userdoc/}{\textbf{FSLeyes}} - an interactive OpenGL-based
desktop application for visualisation of 3D and 4D MRI data, written
in Python.
\item \href{https://fsl.fmrib.ox.ac.uk/fsl/docs/development/management/index.html}{\textbf{The FSL build system}} - an automated modular
build system for FSL, based on the \href{https://docs.conda.io/}{conda}
environment and package manager, inspired by the \href{https://conda-forge.org/}{conda-forge} package ecosystem, and making extensive use of
GitLab CI and Amazon AWS. As part of this work I also added support
for building FSL on Apple Silicon and Linux aarch64 platforms, and
made substantial improvements to the build processes for our C++/CUDA
applications to achieve compatibilty across a wide range of GPU hardware.
\item \href{https://git.fmrib.ox.ac.uk/fsl/MMORF}{\textbf{MMORF}} - a CUDA-accelerated C++ application
for non-linear registration of MRI images, designed and developed by
my colleague \href{https://www.ndcn.ox.ac.uk/team/frederik-lange}{Frederik Lange}. I ported the
codebase to pure C++, using the \texttt{std::thread} library for
parallelism, in order to allow MMORF to be used on a wider range of
platforms, and to take advantage of lower cost cloud-based compute
instances when running at scale.
\item \href{https://github.com/pauldmccarthy/indexed_gzip}{\texttt{indexed\_gzip}} - a Python library for fast
random access of gzip files, written in C and Cython.
\item \href{https://git.fmrib.ox.ac.uk/fsl/armawrap}{\texttt{armawrap}} - a C++ header-only library for
supporting legacy code that was written against the \href{https://www.robertnz.net/nm_intro.htm}{Newmat} linear algebra library but which, under the hood, uses
the more modern \href{https://arma.sourceforge.net}{Armadillo} linear algebra
library.
\item \href{https://pages.fmrib.ox.ac.uk/fsl/funpack/}{\textbf{FUNPACK}}, a Python-based command-line
application for working with very large CSV files from the \href{https://www.ukbiobank.ac.uk/}{UK Biobank}.
\item \href{https://git.fmrib.ox.ac.uk/fsl/pyfeeds}{\texttt{pyfeeds}}, a bespoke testing framework used
to run our internal suite of integration tests, which supports custom
evaluation rules, and can execute tests in parallel on a HPC cluster.
\end{itemize}\leavevmode\newline

Most of the software I develop is open-source, and I actively use and
contribute to open-source projects:\newline

\begin{itemize}
\item During development of my own projects, I have made many small
contributions to a number of open source projects, such as \href{https://github.com/nipy/nibabel/pulls?q=is%3Apr+author%3Apauldmccarthy}{\texttt{nibabel}}, \href{https://github.com/wxWidgets/Phoenix/pulls?q=is%3Apr+author%3Apauldmccarthy}{\texttt{wxPython}},
\href{https://github.com/ipython/ipython/pulls?q=is%3Apr+author%3Apauldmccarthy}{\texttt{IPython}}, \href{https://github.com/pandas-dev/pandas/pull/40180}{\texttt{pandas}},
and \href{https://github.com/mikedh/trimesh/pulls?q=is%3Apr+pauldmccarthy}{\texttt{trimesh}}.
\item I have played an active role in the \href{https://conda-forge.org/}{conda-forge} package ecosystem, having used conda-forge to
publish my own software, and being added to the maintainence team for
a number of other packages, such as \href{https://github.com/conda-forge/wxwidgets-feedstock}{\texttt{wxwidgets}}, \href{https://github.com/conda-forge/wxpython-feedstock}{\texttt{wxPython}}, and \href{https://github.com/conda-forge/libxmlpp-feedstock}{\texttt{libxml++}}.
\end{itemize}\leavevmode\newline

I have also played an active role in my workplace in teaching,
public engagement, and generally helping out where I can:\newline

\begin{itemize}
\item In 2023 I organised and led an effort to migrate the FSL documentation
from an old Wiki framework (\href{https://github.com/moinwiki}{MoinMoin}) to a
markdown-based framework (\href{https://www.mkdocs.org/}{MkDocs}). I organised
a \href{https://git.fmrib.ox.ac.uk/paulmc/fsl-docathlon-2023/}{"docathlon"} to solicit help from colleagues
in porting the old documentation over to the \href{https://fsl.fmrib.ox.ac.uk/fsl/docs/}{new site}, which also doubled as a \texttt{git} training workshop for the
participants.
\item I have volunteered for a range of public engagement activities; for
example, for the past two years I have ran an \href{https://pages.fmrib.ox.ac.uk/paulmc/oxford_neuroscience_experience/}{"MRI analysis bootcamp"} for high-school students interested in
pursuing scientific studies at University.
\item When my boss \href{https://www.nature.com/articles/d41586-023-01391-5}{led the resignation of the NeuroImage
editorial board in 2023}, I did the grunt work behind the scenes to
set up an interim journal submission system (based on \href{https://github.com/BirkbeckCTP/janeway/}{Janeway}) for the new \href{https://direct.mit.edu/imag}{Imaging Neuroscience}
journal. This ran successfully for several months until MIT Press took
over management.
\item I have taught on, and developed large portions of the materials for our
internal \href{https://git.fmrib.ox.ac.uk/fsl/win-pytreat}{"PyTreat"}, a series of internal
Python-based hackathon events, intended to help researchers in getting
started with the Python programming language.
\item I have held lecturing and tutoring roles on our internal \href{https://www.ndcn.ox.ac.uk/study-with-us/graduate-students/fmrib-graduate-programme}{MRI Graduate Programme}, held annually for the
benefit of Oxford-based researchers and PhD students.
\item I have lectured and tutored on the \href{https://fsl.fmrib.ox.ac.uk/fslcourse/}{FSL course}, an international course held annually, which teaches MRI analysis
using FSL.
\end{itemize}\leavevmode\newline

I have been fortunate enough to be listed as a co-author for my
technical contributions to a range of academic studies, including:\newline

\begin{itemize}
\item {\small\small\textbf{Go Figure: Transparency in neuroscience images preserves context and
clarifies interpretation
}\\ \textcolor{grey}{\textit{Taylor et. al. Arxiv [Preprint] 2025}}}\\ \mbox{\small\url{https://doi.org/10.48550/arXiv.2504.07824}}
\item {\small\small\textbf{MMORF—FSL's MultiMOdal Registration Framework}\\ \textcolor{grey}{\textit{Lange et. al. Imaging Neuroscience 2024}}}\\ \mbox{\small\url{http://doi.org/10.1162/imag_a_00100}}
\item {\small\small\textbf{The mouse motor system contains multiple premotor areas and partially follows human organizational principles}\\ \textcolor{grey}{\textit{Lazari et. al. Cell Reports 2024}}}\\ \mbox{\small\url{https://doi.org/10.1016/j.celrep.2024.114191}}
\item {\small\small\textbf{The effects of genetic and modifiable risk factors on brain regions vulnerable to ageing and disease}\\ \textcolor{grey}{\textit{Manuello et. al. Nature Communications 2024}}}\\ \mbox{\small\url{http://doi.org/10.1038/s41467-024-46344-2}}
\item {\small\small\textbf{SARS-CoV-2 is associated with changes in brain structure in UK Biobank}\\ \textcolor{grey}{\textit{Douaud et. al. Nature 2022}}}\\ \mbox{\small\url{https://doi.org/10.1038/s41586-022-04569-5}}
\item {\small\small\textbf{Image processing and Quality Control for the first 10,000 brain imaging datasets from UK Biobank}\\ \textcolor{grey}{\textit{Alfaro-Almagro et. al. Neuroimage 2018}}}\\ \mbox{\small\url{https://doi.org/10.1016/j.neuroimage.2017.10.034}}
\end{itemize}

\end{category}




\begin{category}{2010-2013 - PhD}
\citemnobullet
\begin{changemargin}{-0.5in}{0.5in}
\textcolor{grey}{\textit{Department of Computer Science, University of Otago, Dunedin, New Zealand}}
\end{changemargin}

\citemnobullet My \href{https://ourarchive.otago.ac.nz/handle/10523/4864}{PhD} revolved around the application of graph
theory techinques to fMRI data, to explore how functional
connectivity changes with age and the presence of Alzheimer’s
disease. Part of my PhD work was published in an academic journal:\newline

{\small\small\textbf{The age-related posterior-anterior shift as revealed by voxelwise analysis of functional brain networks}\\ \textcolor{grey}{\textit{McCarthy et. al. Frontiers in Aging neuroscience 2014}}}\\ \mbox{\small\url{https://doi.org/10.3389/fnagi.2014.00301}}\leavevmode\newline

Most of the code which formed the basis for my analyses was written as
part of a package I developed called \href{https://github.com/pauldmccarthy/ccnet}{\texttt{ccnet}},
written in C.\newline

Before migrating my code to C, I used the first 12 months of my PhD as
an opportunity to teach myself Haskell, but then decided to port all
of my code to C for more fine-grained control over performance and
memory usage.

\end{category}




\begin{category}{2007-2010 - Software engineer}
\citemnobullet
\begin{changemargin}{-0.5in}{0.5in}
\textcolor{grey}{\textit{CSIRO, Hobart, Australia}}
\end{changemargin}

\citemnobullet I worked as the primary software engineer in the deployment of three
wireless sensor networks in Tasmania, including:\newline

\begin{itemize}
\item a relatively large scale network (70 nodes) deployed for \href{https://dl.acm.org/doi/10.1145/1435473.1435487}{real time monitoring of soil moisture} in a farm
environment;
\item a \href{https://doi.org/10.1109/OCEANSE.2009.5278177}{multi-hop underwater sensor network}, using
acoustic modems for inter-node communication, and a 3G modem for
internet uplink
\item and a prototype low cost marine wireless sensor node, for deployment in
salmon hatcheries.
\end{itemize}\leavevmode\newline

I played leading roles in each of these projects, working on hardware/
electronic design and construction, low level system programming and
routing protocols in C, and data transmission, storage, and real time
data presentation via a dedicated website, using Java.\newline

As part of another project, I worked on the development of a
platform-agnostic operating system written primarily in C, for low
powered wireless devices with 8-bit microcontrollers. I worked on the
design and development of system and device drivers, including work on
a bootloader with wireless capabilities, for remote reprogramming.\newline

Before leaving CSIRO to begin my PhD, I worked as the lead developer
on a MATLAB toolbox for the processing of marine moorings data, with
the aim of providing a standard output format for data to be hosted by
the Integrated Marine Observing System (IMOS). I am quite proud of the
fact that the project is still actively used and developed - it can be
found at \mbox{\small\url{https://github.com/aodn/imos-toolbox/}}

\end{category}




\begin{category}{2006-2007 - Software engineer}
\citemnobullet
\begin{changemargin}{-0.5in}{0.5in}
\textcolor{grey}{\textit{Australian Defence Force Academy, Canberra, Australia}}
\end{changemargin}

\citemnobullet A mixture of software development, maintenance, and experimentation.
My first project involved the development of a multi-agent system for
crowd behaviour simulation during panic situations. This entailed the
implementation (in Java) of the Helbing particle physics model, which
is a prevalent crowd simulation model. In another project, I developed
a dynamic social network simulation program, also in Java, which
provided visualisation of social network simulations, and exporting in
graph or tabular form of simulation data.

\end{category}




\begin{category}{2002-2006 - Master of Science}
\citemnobullet
\begin{changemargin}{-0.5in}{0.5in}
\textcolor{grey}{\textit{Department of Computer Science, University of South Australia, Adelaide, Australia}}
\end{changemargin}

\citemnobullet I completed a Bachelor of Computer Science with First Class Honours
(GPA 6.5), and used credit from my Honours degree to progress to a
Masters degree. My \href{forensic_analysis_of_mobile_phones.pdf}{Masters thesis} involved an
exploration into different methods of acquiring data from mobile
phones, and the issues involved in presenting this data as forensic
evidence.

\end{category}




\begin{category}{Outside of work}
\citemnobullet
\begin{changemargin}{-0.5in}{0.5in}
\textcolor{grey}{\textit{Llanberis, North Wales, UK}}
\end{changemargin}

\citemnobullet Most of my free time is spent in the outdoors or at the climbing
wall. I have spent much of my adult life suffering from an addiction
to rock climbing and, since moving to North Wales in 2018, have also
developed an affinity for fell running and scrambling. I also enjoy
mountain biking, surfing, and snowboarding when conditions and time
allow. I occasionally post photos and trip reports at \mbox{\small\url{https://pauldmccarthy.github.io/}}.\newline

I am a keen but mediocre guitarist, and also enjoy board games, chess,
reading novels, and \textit{very} occasionally hacking on hobby
coding projects that usually lead nowhere.

\end{category}




\begin{category}{References}
\citemnobullet I am happy to provide references on request.
\end{category}

\end{document}